\section{Growing cutoffs and explicit asymptotic time windows}
\label{sec:growingbox}
We keep the states, projected Gaussian operators, exact-energy convention,
proper group and strict-$G_N\to0$ assumptions of Section~16. The issue is
the dependence of its compact constants on the offsets. We give deliberately
coarse analytic bounds, sufficient for explicit asymptotic windows. They
do not provide a certified numerical minimum mass for a prescribed accuracy.
The proofs below are author self-reviewed drafts, using the Euler and
Stirling estimates of DLMF~\cite[5.11, 15.6, 15.8, 2.3]{DLMF} and the
exact KL kernel~\cite{KL}; no priority claim is made.

\subsection{Constants in the Euler endpoint argument}
Let $Q_0\ge1$ bound the absolute values of the four real representation
offsets, with fixed $\nu>0$. The symbol $Q_0$ is unrelated to either
projection $Q_M$ or the Gaussian quadratic form $\mathcal Q$.
All the small parameters in Sections~10--11 then have imaginary parts
bounded by $c_0(1+Q_0)$, where $c_0$ is independent of $M,Q_0$.
Fix $0<\epsilon<1/4$. A complex tube of radius $\rho_\epsilon>0$
can be chosen independently of $Q_0$ so that every Euler endpoint real
part is between $1/4$ and $3/4$, the Gauss denominator real parts stay
near one, and $|\operatorname{Re}\delta|<\epsilon/4$.
The balanced identity remains exact. We take smaller nested disks when
removing the apparent $e=0$ and $\delta=0$ poles.

We will use an envelope
\begin{equation}
 \mathcal K(Q_0)=C e^{c Q_0}(1+Q_0)^d.
 \label{eq:gb-envelope}
\end{equation}
The fixed constants may be enlarged a finite number of times below;
they depend on $\nu,\epsilon$ and the tube margins, not on $Q_0$ or $M$.
Products and finite sums of such envelopes have the same form.

\begin{lemma}[Parameter-dependent endpoint bounds]
For the four KL hypergeometric factors in Section~11, the endpoint
remainder in \eqref{eq:endpoint-two-branch}, away from the small pole
disks, is bounded by
\begin{equation}
 \mathcal K(Q_0) M^{-1+\epsilon/4},
 \qquad M\ge C_0(1+Q_0)^2.
 \label{eq:gb-endpoint-bound}
\end{equation}
The same envelope class controls the required finite derivative norms
of the endpoint amplitudes. The constants and the large-mass threshold
are uniform over the indicated growing offset boxes.
\label{lem:gb-endpoint}
\end{lemma}
\begin{proof}
There are four separate estimates, with fixed differentiation order.

First, Gamma functions and their reciprocals on fixed vertical strips,
away from the finitely many poles in those strips, are bounded by an
exponential in the imaginary part times a polynomial. Stirling's formula
on the tails and compactness on the bounded part give this bound; fixed
derivatives obey it as well. Thus every small-parameter Gamma product,
connection coefficient $G_\pm$, and $c$ ratio off the pole disks is
bounded by \eqref{eq:gb-envelope}. The pole distance is fixed by
$\rho_\epsilon$, not by $Q_0$.

Second, for $Z=\pm i\lambda$, $\lambda$ proportional to $M$, and shifts
$|a|,|b|\le c_1(1+Q_0)$ with bounded real parts, the sectorial Stirling
remainder gives
\begin{align}
 \frac{\Gamma(Z+a)}{\Gamma(Z+b)}Z^{b-a}&=e^{E(Z;a,b)},\notag\\
 |E(Z;a,b)|&\le C_1(1+Q_0)^2/M,
 \qquad M\ge C_0(1+Q_0)^2.
 \label{eq:gb-gamma-ratio}
\end{align}
Here $E$ is defined using the analytic log-Gamma branch in the fixed
sector about $\pm i\mathbb R_+$, rather than the principal logarithm
of the final ratio. To see the bound directly, write
\[
 E=\int_b^a[\psi(Z+z)-\log Z]\dd z.
\]
On the joining segment $|Z+z|\asymp M$ and its argument stays away
from the negative real axis. The digamma expansion with sectorial
remainder bounds its integrand by $C(1+Q_0)/M$; the segment length is
$O(1+Q_0)$. Exponentiating gives a relative error bounded by
$C(1+Q_0)^2/M$ after increasing $C_0$. This derives a growing-shift
estimate, rather than substituting growing parameters into a fixed-shift
ratio formula. Factors such as $Z^{i y}$ cost at most $e^{C Q_0}$.

Third, all Gauss functions needed near either endpoint have an Euler
integral with real positive bases. For $0\le z\le z_*<1$,
$\operatorname{Re}B,\operatorname{Re}(C-B)>0$, and fixed $j\le5$,
the ordinary Gauss function obeys
\begin{align}
 |\partial_z^j{}_2F_1(A,B;C;z)|
 &\le |(A)_j|\left|\frac{\Gamma(C)}{\Gamma(B)\Gamma(C-B)}\right|
 \mathrm B(\operatorname{Re}B+j,\operatorname{Re}(C-B))\notag\\
 &\hspace{12mm}\times
 (1-z_*)^{-\max(\operatorname{Re}A+j,0)}.
 \label{eq:gb-euler-derivative}
\end{align}
This follows by differentiating Euler's integral: $t,1-t,1-zt$ are
positive, so their imaginary exponents have modulus one. It uses the
ordinary function, with the $\Gamma(C)$ restored from DLMF 15.6.1.
On the tube, the needed $(A,B;C)$ are $(b,d;g)$,
$(b,d;1-\delta)$ and $(g-b,g-d;1+\delta)$; all satisfy the stated
real-part conditions with fixed margins. Their derivative bounds have
form \eqref{eq:gb-envelope}. Near $w=0$ apply the latter two to
$z=1-e^{-w}$ on $0\le w\le1$; for $w$ away from zero apply the first
to $z=e^{-w}$.

Fourth, in \eqref{eq:endpoint-J}, $h(z)=(1-e^{-z})/z$ is positive
on the real half-line, $h(z)\asymp(1+z)^{-1}$, and its logarithmic
derivatives of any fixed positive order are bounded there. Direct
differentiation of $J$ therefore gives, for $j\le5$,
\begin{equation}
 |J^{(j)}(w)|\le C_j(1+Q_0)^{d_j}(1+w)^{d_j}e^{-w/4}.
 \label{eq:gb-J-bound}
\end{equation}
The $t$ weights remain integrable with fixed endpoint margins.
For $w\le1$, the same bounds apply to the smooth factors multiplying
$w^{a-1}$ and $w^0$ in the connection expansion, with
$a=1-\delta$. Their first Taylor remainders have respective forms
$w^a R_1(w)$ and $wR_2(w)$, with derivatives through order four bounded by
$\mathcal K(Q_0)$. Choose a smooth cutoff at zero once and for all,
independent of the parameters. Splitting at $w=1/|Z|$ and integrating
by parts four times on dyadic intervals gives transform bounds
$\mathcal K(Q_0)|Z|^{-\operatorname{Re}a-1}$ and
$\mathcal K(Q_0)|Z|^{-2}$. Differentiating $w^a$ only introduces
fixed-degree polynomials in $a$; the dyadic sums are uniform because
$\operatorname{Re}a$ stays near one. The exponentially decreasing
part away from zero has four integrable derivatives, with the same
envelope. The leading transforms are the Abel boundary values
$\Gamma(a)Z^{-a}$ and $Z^{-1}$, with their cutoff errors handled by
the same integrations by parts.

Finally, \eqref{eq:gb-gamma-ratio} controls the normalization
$\mathcal N_M\mathrm B(\alpha,\beta)$ in Section~10. Multiplying
the transform estimates by it, including its relative error, gives
\eqref{eq:gb-endpoint-bound}. No derivative order grows with $Q_0$.
\end{proof}

\subsection{Assembly and the growth of the compact constants}
\begin{proposition}[Growing-box density control]
There are $C,c,d,C_0>0$, for each fixed $\nu,\sigma,\epsilon$, such
that the compact constants in \eqref{eq:tb-cutoff-bound} may be chosen as
\begin{equation}
 C'_{L,R,\epsilon}+C''_{L,R}
 \le C e^{c(L+R)}(1+L+R)^d,
 \qquad M\ge C_0(1+L+R)^2.
 \label{eq:gb-compact-growth}
\end{equation}
The estimate retains the hard moving threshold. For $F$ one also
requires $M-R>m_*$, which holds eventually for the choices below.
\label{prop:gb-growth}
\end{proposition}
\begin{proof}
For representation offsets bounded by $Q_0$, the exact beta-like
prefactors have relative Gamma-ratio error
$\mathcal K(Q_0)/M$, and size
$\mathcal K(Q_0)M^{-1+\epsilon/4}$ off the apparent poles.
Lemma~\ref{lem:gb-endpoint}, the exact cancellation (11.6), and the
leading coefficients (11.7) consequently give
\begin{equation}
 |S-S_M^0|\le\mathcal K(Q_0)M^{-2+\epsilon}
 \label{eq:gb-raw-bound}
\end{equation}
there, leaving room in $\epsilon$ for the fixed tube margins.
The exact triple-integral holomorphy argument of Section~11 uses
real parts of exponents only. Those margins are independent of the
real offsets; hence its complex tube need not shrink with $Q_0$.
Use fixed-radius disks in $e,\delta$ at their intersections. Their
boundary stays in a box enlarged by a constant, on which the envelope
has the same form. Maximum modulus transfers \eqref{eq:gb-raw-bound}
inside. A pointwise choice of these disks has uniformly bounded overlap;
one does not multiply the bound by an uncontrolled covering constant.
Other Gamma arguments stay near real part $1/2$, so no new poles are
encountered on enlarging the real offsets.

The eight large external Gamma factors satisfy the same growing-shift
Stirling estimates. Their exponentially large mass factors cancel
the hyperbolic prefactor exactly at leading order, as in (11.14).
The remaining small Gamma factors and fixed offset derivatives have
envelope \eqref{eq:gb-envelope}. Thus $F_M/M-H_\nu$ is bounded by
$\mathcal K(Q_0)/M$; the same type of bound applies to the one-step
positive-kernel error. Hyperbolic errors are exponentially small for
$M\ge C_0(1+Q_0)^2$ and are absorbed in these estimates.

In exact-energy coordinates $|\xi|\le L+R$, Taylor expansion gives
\[
 s(M+\xi)=M+\xi-\frac1{8M}
                   +O((1+L+R)M^{-2}).
\]
The Jacobians have uniform $O(M^{-2})$ error when $M>2(L+R+1)$.
After removal of the pole, the first derivatives of the leading bracket
are bounded by $\mathcal K(Q_0)(1+\log(2M))^2$; this follows either
by its explicit sinc subtraction or the small-Gamma bounds above.
The argument of Lemma~16.2 now proves its density bound with an
envelope of the stated form, taking $Q_0=1+L+R$ and enlarging the
constants for fixed $\nu$. Compact volume factors add only a polynomial.
The external density contributes its mass, at most one, rather than
its supremum or derivatives. The Gaussian has modulus at most one.
For the ordered norms, apply the same positive one-step estimate twice.
No growing time seminorm or threshold smoothing constant is introduced.
\end{proof}

\subsection{An explicit cutoff choice for exponential-moment packets}
We now assume a fixed positive exponential moment:
\begin{equation}
 B_0:=\int_0^\infty e^{\beta_0r}\pi_0(r)\dd r<\infty,
 \qquad \beta_0>0.
 \label{eq:gb-moment}
\end{equation}
All fixed Gamma and shifted-exponential packets used above satisfy this,
for example with $\beta_0=\beta/2$; their $B_0$ are respectively
$2^{p+1}$ and $2e^{\beta g/2}$. This extra assumption is not made for
the arbitrary $L^1$ packets of Theorem~16.3.

\begin{theorem}[Rate and asymptotic cutoffs]
For a packet satisfying \eqref{eq:gb-moment}, both protocols obey,
for every fixed $q>0$,
\begin{equation}
 \delta_M^P:=\sup_{t\in\mathbb R}|D_M^P(t)-D_M^{(2),P}(t)|
 =O_q\bigl(\ell^{-q}\bigr),\qquad \ell=\log(2M).
 \label{eq:gb-log-rate}
\end{equation}
The same statement holds separately for the crossed error and the norm
error. An explicit admissible cutoff choice, for sufficiently large $M$, is
\begin{equation}
 L_M=\frac{q+2}{\beta_0}\log\ell,
 \qquad R_M=\sqrt{\frac{8(q+2)}{\sigma^2}\log\ell}.
 \label{eq:gb-explicit-cutoffs}
\end{equation}
There is also a positive mass-power rate $\delta_M^P=O(M^{-\kappa})$
for some $\kappa>0$ depending on the envelope constants and the packet.
We do not evaluate a numerical value of $\kappa$ here.
\label{thm:gb-rate}
\end{theorem}
\begin{proof}
Markov's inequality gives $\Pr(r>L)\le B_0e^{-\beta_0L}$.
For $R\ge1$, the Gaussian envelope of Section~9 gives
\[
 b(R/2)\le C e^{-\sigma^2R^2/4},\quad
 \tau_X+\tau_N\le C\{e^{-\sigma^2R^2/8}+e^{-\beta_0L}\}.
\]
With \eqref{eq:gb-explicit-cutoffs}, both tails are
$O(\ell^{-q-2})$. Write $Q_0=1+L_M+R_M=O_q(\log\ell)$.
The compact term in Theorem~16.3 and
Proposition~\ref{prop:gb-growth} is then at most
\[
 C_q\ell^{a_q+o(1)}
 \left[M^{-1+\epsilon}+\frac{1+\ell}{M}
                        +\frac{(1+\ell)^2}{M^2}\right],
\]
for a finite constant $a_q$. For each fixed $q$ this is smaller than
every inverse power of $\ell$, and $M\ge C_0Q_0^2$ holds eventually.
This proves \eqref{eq:gb-log-rate}, with no exchange of a fixed-box
limit and a growing box.

For the last assertion set $\epsilon=1/8$, let $c$ be the enlarged
coefficient in \eqref{eq:gb-compact-growth}, and put
\[
 \kappa=\frac{\beta_0}{8(c+\beta_0)},\quad
 L_M=\frac{\log M}{4(c+\beta_0)},\quad
 R_M=\sqrt{\frac{16\kappa}{\sigma^2}\log M}.
\]
Both tails are $O(M^{-2\kappa})$. The compact term is
$M^{-7/8+c/[4(c+\beta_0)]+o(1)}$, which is $O(M^{-\kappa})$
since $\kappa\le1/8$. The other terms are smaller; the required
large-mass condition still holds. The constants $c,C_0$ have analytic
existence bounds above but have not been numerically certified.
\end{proof}

\subsection{Consequences for the exact finite-mass observable}
Let $c_p\ne0$ be the real coefficient in \eqref{eq:dep-real-tail},
and $n=p+3$ for even $p$, $n=p+4$ for odd $p$. All density, field,
Gaussian and projection parameters are fixed in the following limits.
For every fixed $0<\alpha<1$, the \emph{exact} commutator satisfies
\begin{equation}
 D_M^p(\ell-\ell^\alpha)-\mathcal D_{{\rm edge},p}
 =c_p\ell^{-\alpha n}
  \left[1+O\left(\ell^{-\alpha}
                         +\ell^{-n(1-\alpha)}\right)\right].
 \label{eq:gb-gamma-window}
\end{equation}
In particular $s=\sqrt\ell$, which was not justified by Section~16
alone, is now an explicit asymptotic window. The signed coefficient
remains negative for $p=1,2$, $\beta=1$.
For a fixed fraction $0<a<2$ one instead obtains
\begin{equation}
 D_M^p((1-a)\ell)-\mathcal D_{{\rm edge},p}
 =c_p\ell^{-n}\{a^{-n}+(2-a)^{-n}\}
                         +O(\ell^{-n-1}).
 \label{eq:gb-gamma-fraction}
\end{equation}
This estimate is uniform on closed subintervals of $0<a<2$.
It retains both tails, including for positive physical time $a<1$.
For fixed physical $t$ in a bounded interval, it implies
\begin{equation}
 D_M^p(t)-\mathcal D_{{\rm edge},p}
 =2c_p\ell^{-n}+O(\ell^{-n-1}).
 \label{eq:gb-fixed-time-rate}
\end{equation}
For bounded $T$ at the positive logarithmic-time transition, it also gives
\begin{equation}
 D_M^p(\ell+T)-\mathcal D_p(T)
 =c_p(2\ell+T)^{-n}+O(\ell^{-n-1}).
 \label{eq:gb-fixed-T-rate}
\end{equation}
The critical odd-$p$ case with $c_p=0$ is excluded from these leading
coefficient statements.

For $F$, every fixed $b>0$ and $2<h<4$ give a different explicit window:
\begin{equation}
 D_M^F(\ell-b\log\ell)
 =2C_{\nu,\sigma}\ell^{-2b}
                  [1+O_h(\ell^{-(h-2)b})].
 \label{eq:gb-F-window}
\end{equation}
These formulas follow by inserting the Section~13--15 tails in the
mirror identity and choosing $q$ in \eqref{eq:gb-log-rate} larger than
the power being resolved. For \eqref{eq:gb-gamma-window}, the fast tail
relative to the slow one is $O(\ell^{-n(1-\alpha)})$; for
\eqref{eq:gb-F-window} it is exponentially small in $\ell$.
No unknown $\delta_M$ appears in the time choices.

For a fixed gap $g>0$, \eqref{eq:tb-gap-window} transfers to the exact
$D_M^g$ with its displayed $O(s^{-5}+k^{-5})$ error whenever $s,k$
are fixed positive powers of $\ell$ or fixed fractions away from zero.
This is an absolute envelope result. A relative statement still requires
separation from zeros of the relevant oscillatory combination, and the
estimate is not uniform as $g\downarrow0$.

The asymptotic window shapes are explicit, but these theorems do not
give a numerical mass at which a chosen percentage accuracy begins.
They are restricted to the stated free-matter, $G_N\to0$ model and do
not identify a finite-gravity chaos exponent. In particular the
$M^{-4}$ mirror correction for $F$ at fixed $T$ need not be the leading
exact correction; the above error bound need not resolve that power.

\subsection{Diagnostic scope and next physical comparison}
The new diagnostics check the shifted Gamma ratios on growing-parameter
fixtures, the explicit Euler derivative inequality and exact-energy
coordinate estimates. These probe the ingredients and conventions of
the proof. They do not certify its supremum constants, perform a full
finite-$M$ 6j packet integration, or independently validate the proof.
Previously frozen limiting spectra are not recomputed.

The saved run contains 290 checks at 80 decimal digits, with 110-digit
Gamma-ratio cross-checks. On the recorded shifts up to $Q_0=64$, the
largest $M|E|/(1+Q_0)^2$ is about $0.080004$; this is a sampled
quantity, not the universal constant in \eqref{eq:gb-gamma-ratio}.
The largest sampled ratio of a Gauss derivative to its Euler-integral
bound is about $0.883322$. The energy-coordinate and normalized tail
inequalities are checked separately. No finite collection of these
values proves the growing-box supremum estimate.

With the explicit asymptotic window established at this level, the next
physical step is a matched finite-$G_N$ comparison. The external state,
projection, proper time, normalization, matter representation and contour
must agree before the current preparation effects can be compared with
a genuine gravitational interaction or a microscopic candidate.
Numerically useful error constants and exact finite-mass packet checks
remain a separate validation task.
